This study presents an explainable AI-based decision support system for precision allocation of agricultural skill development programs. By integrating segmentation, predictive modeling, and explainable analysis into a unified framework, the proposed approach demonstrates how artificial intelligence can be applied to support decision-making in human capital development within agricultural systems.

The results show that farmer populations exhibit significant heterogeneity, and that key drivers of skill improvement are associated with baseline capability, risk perception, and institutional support. The application of SHAP analysis further reveals non-linear and segment-specific effects, providing deeper insights beyond conventional analytical methods.

A central contribution of this work lies in transforming explainable AI outputs into explicit and interpretable decision rules. This enables model insights to be operationalized in practice, supporting practitioners in making informed and transparent training allocation decisions without requiring direct interaction with machine learning models.

The findings highlight that the role of artificial intelligence in agriculture extends beyond predictive accuracy toward the development of deployable and interpretable decision systems. In this context, the proposed Segment--Attribute--Recommend framework provides a structured and practical approach for bridging the gap between data analytics and real-world implementation.

Despite these contributions, the study is subject to several limitations. The analysis is based on observational data, and results are interpreted as associative rather than causal. In addition, the framework is developed within a specific empirical context, and further validation is needed to assess its generalizability across different agricultural systems.

Future research should explore the integration of dynamic data sources, real-time decision support tools, and prospective validation through experimental or quasi-experimental designs. Such extensions would further enhance the ability of AI-driven systems to support adaptive and scalable agricultural development strategies.
